\chapter{Bangun Datar dan Letak}\label{ch:g1:shapes}

Jendela berbentuk persegi panjang, roda berbentuk lingkaran, dan atap
menyembunyikan segitiga: bangun datar ada di mana-mana. Bab ini
melatih kita mengenalinya dengan membilang sisi dan titik sudut, lalu
mengatakan dengan tepat \emph{di mana} letak sesuatu: kiri, kanan, di
atas, di bawah, di antara.

\section{Empat bangun datar dasar}

\begin{definition}[Persegi, persegi panjang, segitiga, lingkaran]\label{def:g1:shapes:def}
\begin{itemize}
  \item \emph{segitiga}\index{segitiga} punya $3$ sisi dan $3$ titik sudut;
  \item \emph{persegi}\index{persegi} punya $4$ sisi sama panjang dan $4$ titik sudut;
  \item \emph{persegi panjang}\index{persegi panjang} punya $4$ sisi dan $4$ titik
        sudut, dengan sisi berhadapan yang sama panjang (persegi yang
        ditarik memanjang);
  \item \emph{lingkaran}\index{lingkaran} bundar sempurna: tanpa sisi, tanpa titik sudut.
\end{itemize}
\end{definition}

\begin{omfigure}
\begin{tikzpicture}[scale=0.7]
  \draw[thick, fill=omDef!12] (0,0) -- (2,0) -- (1,1.7) -- cycle;
  \node[below, align=center, font=\small] at (1,-0.25)
    {segitiga:\\$3$ sisi};
  \begin{scope}[xshift=3.6cm]
  \draw[thick, fill=omThm!12] (0,0) rectangle (1.7,1.7);
  \node[below, align=center, font=\small] at (0.85,-0.25)
    {persegi:\\$4$ sisi sama panjang};
  \end{scope}
  \begin{scope}[xshift=6.9cm]
  \draw[thick, fill=omMeth!12] (0,0.05) rectangle (2.8,1.65);
  \node[below, align=center, font=\small] at (1.4,-0.25)
    {persegi panjang:\\$4$ sisi};
  \end{scope}
  \begin{scope}[xshift=11.3cm]
  \draw[thick, fill=omProp!12] (0.9,0.85) circle (0.9);
  \node[below, align=center, font=\small] at (0.9,-0.25)
    {lingkaran:\\tanpa titik sudut};
  \end{scope}
\end{tikzpicture}

{\small Empat bangun datar dasar. Untuk menamai sebuah bangun, jangan
percaya pada pandangan pertama: \emph{bilanglah} sisi dan titik
sudutnya.}
\end{omfigure}

\begin{example}[Bangun yang miring tetap bangun yang sama]\label{ex:g1:shapes:tilted}
Persegi yang berdiri di atas titik sudutnya tetaplah persegi ---
empat sisi sama panjang, empat titik sudut --- walaupun terlihat
seperti layang-layang. Memutar sebuah bangun tidak mengubah jenisnya.
\end{example}

\begin{example}[Berburu bangun datar]\label{ex:g1:shapes:hunt}
Di dalam kelas: papan tulis berbentuk persegi panjang, jam dinding
berbentuk lingkaran, serbet yang dilipat bisa menjadi segitiga, dan
sebagian jendela berbentuk persegi. Benda nyata bukanlah bangun yang
sempurna, tetapi permukaannya mendekati bangun itu (bangun ruang dan
sisinya kembali pada \cref{ch:g2:shapes}).
\end{example}

\section{Di mana letaknya?}

\begin{definition}[Kata letak]\label{def:g1:shapes:position}
Untuk mengatakan di mana letak sesuatu: \emph{kiri} / \emph{kanan},
\emph{atas} / \emph{bawah}, \emph{di atas} / \emph{di bawah},
\emph{di depan} / \emph{di belakang}, \emph{di antara},
\emph{di dalam} / \emph{di luar}.
\end{definition}

\begin{omfigure}
\begin{tikzpicture}[scale=0.75]
  \draw[thick, fill=omDef!12] (0,0) circle (0.55);
  \draw[thick, fill=omThm!12] (1.9,-0.55) rectangle (3.0,0.55);
  \draw[thick, fill=omMeth!12] (4.4,-0.5) -- (5.5,-0.5) -- (4.95,0.45)
    -- cycle;
  \node[below, font=\small] at (0,-0.75) {lingkaran};
  \node[below, font=\small] at (2.45,-0.75) {persegi};
  \node[below, font=\small] at (4.95,-0.75) {segitiga};
  \node[right, font=\small, align=left] at (6.3,0)
    {persegi berada \emph{di antara}\\lingkaran dan segitiga;\\
    lingkaran ada di \emph{kiri} persegi};
\end{tikzpicture}

{\small Kata letak sedang bekerja. Hati-hati: kiri \emph{menurutmu}
dan kiri orang yang menghadapmu itu berlawanan!}
\end{omfigure}

\section{Jalur di kertas berpetak}

\begin{method}[Menjelaskan sebuah jalur]\label{met:g1:shapes:path}
Di kertas berpetak, sebuah jalur adalah daftar langkah:
$\to$ (satu petak ke kanan), $\leftarrow$ (ke kiri), $\uparrow$ (ke
atas), $\downarrow$ (ke bawah).
\begin{enumerate}
  \item Untuk mengikuti jalur: mulailah dari petak yang ditandai lalu
        turuti anak panahnya satu per satu;
  \item untuk menjelaskan sebuah jalur: telusuri jalur itu dalam
        kepalamu lalu tulislah anak panahnya secara berurutan;
  \item dua jalur yang berbeda dapat menghubungkan dua petak yang
        sama --- ada yang lebih pendek, ada yang lebih panjang.
\end{enumerate}
\end{method}

\begin{omfigure}
\begin{tikzpicture}[scale=0.6]
  \draw[gray!40] (0,0) grid (7,4);
  \fill[omDef!40] (0.5,0.5) +(-0.42,-0.42) rectangle +(0.42,0.42);
  \node[font=\footnotesize] at (0.5,0.5) {mulai};
  \fill[omThm!40] (5.5,2.5) +(-0.42,-0.42) rectangle +(0.42,0.42);
  \node[font=\footnotesize] at (5.5,2.5) {tujuan};
  % tujuh anak panah satu satuan, satu per langkah, dengan sela kecil
  \foreach \a/\b in {0.5/1.5, 1.5/2.5}
    \draw[-Stealth, omDef, very thick] (\a+0.1,0.5) -- (\b-0.1,0.5);
  \foreach \a/\b in {0.5/1.5, 1.5/2.5}
    \draw[-Stealth, omDef, very thick] (2.5,\a+0.1) -- (2.5,\b-0.1);
  \foreach \a/\b in {2.5/3.5, 3.5/4.5}
    \draw[-Stealth, omDef, very thick] (\a+0.1,2.5) -- (\b-0.1,2.5);
  \draw[-Stealth, omDef, very thick] (4.6,2.5) -- (5.05,2.5);
\end{tikzpicture}

{\small Sebuah jalur dari petak mulai ke petak tujuan, satu anak
panah untuk satu langkah:
$\to\ \to\ \uparrow\ \uparrow\ \to\ \to\ \to$ --- tujuh langkah.
Bisakah kamu menemukan jalur lain dengan \omterm{def:g1:addition:def}{jumlah} langkah yang sama?}
\end{omfigure}

\section{Latihan}

\begin{exercise}[$\star$]\label{exo:g1:shapes:1}
Untuk setiap bangun, bilanglah sisi dan titik sudutnya lalu namailah:
bangun dengan $3$ titik sudut; bangun bundar tanpa titik sudut;
bangun dengan $4$ sisi sama panjang.
\end{exercise}

\begin{exercise}[$\star$]\label{exo:g1:shapes:2}
Carilah di rumah: dua persegi panjang, satu lingkaran, satu segitiga,
satu persegi. (Pintu? piring? rambu jalan?)
\end{exercise}

\begin{exercise}[$\star$]\label{exo:g1:shapes:3}
Benar atau salah? ``Persegi yang diputar sampai berdiri di atas titik
sudutnya berhenti menjadi persegi.'' Jelaskan dengan
\cref{ex:g1:shapes:tilted}.
\end{exercise}

\begin{exercise}[$\star$]\label{exo:g1:shapes:4}
Gambarlah di kertas berpetak: persegi berukuran $3$ petak kali $3$
petak; persegi panjang $5$ kali $2$; segitiga yang memakai diagonal
petak.
\end{exercise}

\begin{exercise}[$\star$]\label{exo:g1:shapes:5}
Letakkan tiga mainan berjajar. Jelaskan barisan itu dengan kata
letak: mana yang berada di antara yang lain? Mana yang di sebelah
kiri?
\end{exercise}

\begin{exercise}[$\star$]\label{exo:g1:shapes:6}
Gambarlah sebuah lingkaran. Gambarlah tanda silang \emph{di dalam}
lingkaran itu, bintang \emph{di luar}, dan titik tepat \emph{pada}
garis lingkarannya.
\end{exercise}

\begin{exercise}[$\star$]\label{exo:g1:shapes:7}
Di kertas berpetak, tandailah petak awal. Ikutilah:
$\to\ \uparrow\ \to\ \uparrow\ \to$. Tandailah petak tibanya. Lalu
tulislah jalur yang kembali lurus ke petak awal.
\end{exercise}

\begin{exercise}[$\star$]\label{exo:g1:shapes:8}
Berapa titik sudut seluruhnya: dua segitiga dan satu persegi? Satu
lingkaran dan dua persegi panjang?
\end{exercise}

\begin{exercise}[$\star\star$]\label{exo:g1:shapes:9}
Ambillah $12$ batang korek api. Bisakah kamu membentuk persegi dengan
memakai \emph{semua} batang itu, dengan batang yang utuh dan \omterm{def:g1:addition:def}{jumlah}
yang sama di setiap sisi? Berapa batang untuk satu sisi? Bisakah
dengan $10$ batang korek api?

\end{exercise}
